\subsection{Matriz de evaluación ponderada}

La selección se formaliza mediante una matriz ponderada basada en los criterios de la Tabla \ref{tab:criterios_compresion}.
Se utiliza una escala discreta de 1 a 5, donde 5 representa el mejor desempeño relativo frente al criterio evaluado.
Los puntajes se asignan con base en la documentación citada y la comparación cualitativa reportada en dichas fuentes.

\begin{table}[H]
    \centering
    \caption{\\Evaluación de alternativas de compresión (escala 1--5) y total ponderado}
    \label{tab:evaluacion_compresores}
    \begin{tabular}{lcccccc}
        \hline
        \multicolumn{1}{l}{\textbf{Criterio}} &
        \multicolumn{1}{c}{\textbf{Peso}} &
        \multicolumn{1}{c}{\textbf{Zstd}} &
        \multicolumn{1}{c}{\textbf{LZ4}} &
        \multicolumn{1}{c}{\textbf{Snappy}} &
        \multicolumn{1}{c}{\textbf{zlib}} \\
        \hline
        C1 (Ratio)                & 0.40 & 4 & 3 & 2 & 5 \\
        C2 (Velocidad)            & 0.20 & 4 & 5 & 4 & 2 \\
        C3 (Overhead memoria)     & 0.15 & 3 & 5 & 5 & 4 \\
        C4 (Paralelismo)          & 0.15 & 5 & 4 & 4 & 3 \\
        C5 (Integración C/C++)    & 0.10 & 5 & 5 & 4 & 5 \\
        \hline
        \textbf{Total}            & 1.00 &
        \textbf{4.10} & \textbf{4.05} & \textbf{3.35} & \textbf{3.95} \\
        \hline
    \end{tabular}
\end{table}

La tabla refleja el compromiso clásico entre ratio y velocidad: zlib tiende a maximizar la compresión, pero penaliza en throughput
(zlib authors, 2022). LZ4 y Snappy destacan por velocidad y bajo overhead, usualmente con ratios inferiores a zlib y Zstd
(lz4 contributors, s.f.; google/snappy contributors, s.f.; Snappy project, s.f.). Zstd se ubica como una alternativa de equilibrio,
combinando ratios altos con rendimiento competitivo y soporte multihilo (facebook/zstd contributors, s.f.; Zstandard project, s.f.).

Como referencia complementaria, una comparación práctica de ratios en datos generales reporta el orden
zlib $>$ Zstandard $>$ LZ4 $>$ Snappy (StarRocks Documentation, s.f.), y benchmarks recientes ilustran el espectro de compromisos
entre ratio, velocidad y memoria (MaskRay, 2025).
